\section{How to use the QPU website}

This chapter assumes no programming experience. You can build circuits by
clicking controls, or edit the text protocol directly. Both paths describe the
same circuit.

\subsection{Five programming words used in this guide}

\begin{description}
  \item[Source text] The human-readable instructions in the protocol editor.
  It is ordinary text, like a note, with one QPU instruction per logical line.
  \item[Instruction] One action for the QPU system, such as SET, H, CNOT, or
  MEASURE. The first word says what action to perform.
  \item[Argument] Extra information after an instruction. In
  \code{SET Q 0p}, Q and 0p are arguments: they say which wire and which state.
  \item[Compile] Translate source text into the ordered visual gates that the
  simulator can run. Compilation checks spelling, required inputs, and targets.
  \item[Run] Apply the compiled gates to the selected starting state. Running
  changes simulator state; it does not rewrite the source text.
\end{description}

\subsection{Interface words and how they relate}

\begin{center}
\begin{tabularx}{\textwidth}{>{\bfseries}l X}
\toprule
Word on screen or in source & Meaning\\
\midrule
Particle / qubit & One simulated quantum two-state system. The interface may
say particle where circuit theory says qubit.\\
Wire & The horizontal visual path followed by one qubit through the circuit.\\
Token & A source-text name, such as Q or Carry, that resolves to a wire.\\
Register & One or more qubit wires considered as a group.\\
Gate & A box or controlled symbol that transforms state.\\
Target & The wire changed by a gate.\\
Control & A wire whose state decides whether a controlled gate acts.\\
Protocol / process & The complete source description of a named circuit.\\
\bottomrule
\end{tabularx}
\end{center}

\subsection{Run a bundled example without writing code}

\begin{enumerate}
  \item Open the navigation menu and choose \textbf{Circuit builder}.
  \item Find \textbf{Load a starter circuit} and choose an example.
  \item Inspect the horizontal wires and gate blocks on the circuit canvas.
  \item Choose \textbf{Play Sequence} to watch gates advance, or
  \textbf{Run all} to jump to the final state.
  \item Choose \textbf{Measure all} if the circuit has not already measured
  its outputs.
  \item Read the particle/state display and runtime log. A measured qubit is a
  classical 0 or 1 outcome for that run.
  \item Choose \textbf{Reset state} before repeating the same experiment.
\end{enumerate}

\subsection{Compile a complete protocol}

Listings marked \textbf{Complete runnable protocol} may be copied as a whole:

\textbf{Complete runnable protocol}
\begin{lstlisting}
MAIN-PROCESS FirstWebsiteCircuit
CREATETOKEN -I Q
SET Q 0p
H -I Q -O Q
MEASURE -I Q
RETURNVALS Q
\end{lstlisting}

\begin{enumerate}
  \item Scroll to \textbf{Compile text into a circuit}.
  \item Replace the editor contents with the complete listing.
  \item Choose \textbf{Compile protocol}.
  \item Read \textbf{Protocol checks}. Zero errors means compilation may
  proceed. Warnings identify accepted syntax that may not behave as expected.
  \item Return to the run controls and choose \textbf{Run all}.
\end{enumerate}

\textbf{Expected result:} the H gate prepares equal 0 and 1 amplitudes.
MEASURE returns either 0 or 1. Repeating from Reset state should produce both
outcomes over many runs.

\subsection{Complete examples versus fragments}

This documentation uses three listing labels:
\begin{description}
  \item[Complete runnable protocol] Copy the entire block into the protocol
  editor and compile it.
  \item[Instruction fragment] A few lines demonstrating one syntax rule. Add
  them to a complete process rather than compiling the fragment alone.
  \item[Companion metadata] Truth-table text saved separately as
  \code{.qpuio}; do not paste it into the protocol editor.
\end{description}

\subsection{When protocol checks report a problem}

An \textbf{error} blocks successful compilation. A \textbf{warning} means the
text is accepted but may use inactive, incomplete, or surprising behavior.
Each diagnostic gives a line number, the relevant source line when available,
and a suggested correction.

Choose \textbf{Open Correction Lab for guided help} below the diagnostic list
when the correction is not obvious. The Correction Lab can test a protocol,
probe outputs, explain failures, and propose gate-level guidance. It does not
change protected bundled truth tables.

\subsection{Saving and reopening work}

Choose \textbf{Download as .qpucir} to save protocol source. The
\code{-qpucir.txt} option contains the same source for devices that only handle
plain text files. Use the navigation menu's \textbf{Upload files} page to
reopen a protocol, optionally together with its matching \code{.qpuio}
truth-table metadata.
